\documentclass[11pt]{article}
\usepackage[margin=1in]{geometry}
\usepackage[T1]{fontenc}
\usepackage[utf8]{inputenc}
\usepackage{amsmath,amssymb,amsthm,mathtools}
\usepackage{graphicx}
\usepackage{xcolor}
\usepackage{microtype}
\usepackage[colorlinks=true,linkcolor=blue!55!black,citecolor=blue!55!black,urlcolor=blue!55!black]{hyperref}

\newtheorem{theorem}{Theorem}
\newtheorem{lemma}{Lemma}
\newcommand{\R}{\mathbb R}

\title{Reflexivity is unnecessary in Ricceri's compact-operator minimax identity}
\author{Candidate full-solution packet for arXiv:1508.04891, Remark 3}
\date{11 August 2026}

\begin{document}
\maketitle

\begin{center}
\fbox{\parbox{0.92\textwidth}{\textbf{Status: candidate full solution, likely valid, subject to human review.}
Theorem 2 of the source remains true on every infinite-dimensional Banach
space.  Thus the counterexample conjectured in a Schur space does not exist.}}
\end{center}

\section{The question and the stronger answer}

Ricceri's Theorem 2 assumes that $E$ is an infinite-dimensional reflexive
real Banach space and proves, for a nonzero compact linear operator $T:E\to E$,
a finite-valued lower-semicontinuous convex coercive functional
$\varphi:E\to\R$, a compact interval $I\ni0$, and a finite-valued
lower-semicontinuous convex $\psi:I\to\R$, that
\begin{equation}\label{eq:source}
 \sup_{x\in X}\inf_{\lambda\in I}
 \bigl(\varphi(Tx-\lambda x)+\psi(\lambda)\bigr)
 =r+\psi(0),
 \qquad X=\{x:\varphi(Tx)\le r\},
\end{equation}
whenever $r>\varphi(0)$ \cite{source}.  Remark 3 asks whether reflexivity can
be dropped and conjectures a counterexample in a Banach space with the Schur
property.

The answer is stronger than requested.

\begin{theorem}[Reflexivity-free identity]\label{thm:main}
Let $E$ be any infinite-dimensional real Banach space, let $T:E\to E$ be a
nonzero compact linear operator, and let $\varphi:E\to\R$ be finite-valued,
lower semicontinuous, convex, and coercive.  Let $I\subset\R$ be any compact
set containing $0$, and let $\psi:I\to\R$ be finite-valued and lower
semicontinuous.  Then \eqref{eq:source} holds for every $r>\varphi(0)$.
\end{theorem}

In particular, neither reflexivity of $E$, convexity of $\psi$, nor the
interval structure of $I$ is needed.  The conclusion therefore holds on
$\ell_1$ and on every other infinite-dimensional Schur space.

\section{Two elementary lemmas}

\begin{lemma}[Compact approximate kernel]\label{lem:kernel}
If $H\subset E$ is an infinite-dimensional closed subspace and $T:E\to E$
is compact, then for every positive sequence $\varepsilon_n\downarrow0$
there are $z_n\in H$ with $\|z_n\|=1$ and $\|Tz_n\|\le\varepsilon_n$.
\end{lemma}

\begin{proof}
Otherwise $T|_H$ is bounded below by some positive constant.  Its inverse on
$T(H)$ is then Lipschitz, while compactness makes the image of the unit ball
of $H$ relatively compact.  The unit ball of $H$ would be relatively compact,
contradicting infinite-dimensionality.
\end{proof}

\begin{lemma}[Intermediate range level]\label{lem:level}
For every $s>\varphi(0)$ there is $y\in E$ such that $\varphi(Ty)=s$.
\end{lemma}

\begin{proof}
Choose $v\ne0$ in $T(E)$.  A finite-valued lower-semicontinuous convex
functional on a Banach space is norm-continuous.  Hence
$t\mapsto\varphi(tv)$ is continuous, equals $\varphi(0)$ at zero, and tends
to $+\infty$ as $t\to+\infty$ by coercivity.  The intermediate value theorem
gives $t\ge0$ with $\varphi(tv)=s$; write $tv=Ty$.
\end{proof}

\section{Proof of the theorem}

\begin{proof}[Proof of Theorem~\ref{thm:main}]
The upper bound is immediate: for $x\in X$, the admissible choice $\lambda=0$
gives
\[
 \inf_{\lambda\in I}\bigl(\varphi(Tx-\lambda x)+\psi(\lambda)\bigr)
 \le \varphi(Tx)+\psi(0)\le r+\psi(0).
\]

Fix an arbitrary $s$ with $\varphi(0)<s<r$.  By Lemma~\ref{lem:level}, choose
$y\in E$ and put
\[
 a=Ty,\qquad \varphi(a)=s.
\]
Continuity and convexity of $\varphi$ give a subgradient $p\in E^*$ at $a$:
\begin{equation}\label{eq:subgradient}
 \varphi(w)\ge \varphi(a)+p(w-a)\qquad(w\in E).
\end{equation}
The closed subspace $H=\ker p$ is infinite-dimensional (also when $p=0$).
Apply Lemma~\ref{lem:kernel} in $H$ to choose
\[
 z_n\in H,\qquad \|z_n\|=1,\qquad \|Tz_n\|\le n^{-2}.
\]
Set $x_n=y+nz_n$.  Then
\[
 Tx_n=a+nTz_n\longrightarrow a,
\]
so $\varphi(Tx_n)\to s<r$.  After discarding finitely many indices,
$x_n\in X$.

For each such $n$, lower semicontinuity on compact $I$ gives a minimizer
$\lambda_n\in I$ and
\[
 m_n:=\min_{\lambda\in I}
 \bigl(\varphi(Tx_n-\lambda x_n)+\psi(\lambda)\bigr)
 =\varphi(Tx_n-\lambda_nx_n)+\psi(\lambda_n).
\]
We claim
\begin{equation}\label{eq:liminf}
 \liminf_{n\to\infty}m_n\ge s+\psi(0).
\end{equation}
It suffices to start with an arbitrary subsequence and pass to a further one
on which $\lambda_n\to\lambda\in I$.

If $\lambda\ne0$, then $Tx_n$ is bounded, while
$\|x_n\|\ge n-\|y\|$; hence
$\|Tx_n-\lambda_nx_n\|\to\infty$.  Coercivity makes the $\varphi$ term tend
to $+\infty$.  Since a finite-valued lower-semicontinuous function on compact
$I$ is bounded below, this case gives more than \eqref{eq:liminf}.

Suppose now that $\lambda=0$ and put $c_n=n\lambda_n$.  If $(c_n)$ is
unbounded, pass further so that $|c_n|\to\infty$.  From
\[
 Tx_n-\lambda_nx_n
 =a+nTz_n-\lambda_ny-c_nz_n
\]
and $n\|Tz_n\|\le n^{-1}$, the residual norm again tends to infinity, so
coercivity applies.

It remains only the case where $(c_n)$ is bounded.  Since $p(z_n)=0$,
\[
 p\bigl(Tx_n-\lambda_nx_n-a\bigr)
 =n\,p(Tz_n)-\lambda_n p(y)\longrightarrow0.
\]
The supporting inequality \eqref{eq:subgradient} therefore gives
\[
 \liminf_n\varphi(Tx_n-\lambda_nx_n)\ge\varphi(a)=s.
\]
Also $\liminf_n\psi(\lambda_n)\ge\psi(0)$ by lower semicontinuity at zero.
This proves \eqref{eq:liminf} in every possible subsequential regime.

Consequently,
\[
 \sup_{x\in X}\inf_{\lambda\in I}
 \bigl(\varphi(Tx-\lambda x)+\psi(\lambda)\bigr)
 \ge \liminf_n m_n\ge s+\psi(0).
\]
Letting $s\uparrow r$ gives the lower bound $r+\psi(0)$, which combines with
the first paragraph to prove equality.
\end{proof}

\section{Why the construction works}

There is one delicate parameter scale.  If a minimizing $\lambda_n$ stays
away from zero, then the term $-\lambda_n nz_n$ escapes to infinity.  The
same is true if $\lambda_n\to0$ but $n|\lambda_n|\to\infty$.  Only
$\lambda_n=O(1/n)$ can keep the residual bounded.  At exactly that scale,
the supporting functional detects all first-order displacement except the
large $z_n$ term, because $z_n$ was chosen in $\ker p$.  Compactness makes
the remaining error $nTz_n$ vanish.

This also explains why reflexive weak compactness in the source proof is not
intrinsic to the identity.  It is one route to a general minimax theorem, but
the particular compact-operator expression admits a direct near-kernel
construction.

\section{Novelty and limitations}

On 11 August 2026, exact searches of the run registry and cheap indexes found
no prior packet or known answer.  Searches across the locally parsed arXiv
corpus found only the source and later bibliographic citations.  A bounded
citation-index query returned two citing works.  Inspection of the primary
texts---Ricceri's 2020 gradient-systems preprint \cite{gradient} and
\emph{Miscellaneous Applications of Certain Minimax Theorems II}
\cite{misc}---found no discussion or resolution of Remark 3.  This is a
bounded novelty check, not an exhaustive priority claim.

The theorem retains the source assumptions that $T$ is nonzero and compact,
$E$ is infinite-dimensional, and $\varphi$ is finite-valued, convex, and
coercive.  The source already notes that infinite-dimensionality and
compactness cannot simply be removed.  The present result addresses only the
reflexivity question in Remark 3; it does not alter the general minimax
Theorem 1.

\appendix
\section{Source-page evidence}

Figure~\ref{fig:source} reproduces the source page containing Theorem 2 and
Remark 3.

\begin{figure}[p]
\centering
\includegraphics[height=0.89\textheight]{figures/open_problem_crop.png}
\caption{Source paper: Theorem 2 and Remark 3, including the conjectured
Schur-space counterexample.}\label{fig:source}
\end{figure}

\begin{thebibliography}{9}
\bibitem{source}
B. Ricceri,
\emph{A minimax theorem in infinite-dimensional topological vector spaces},
arXiv:1508.04891v3 [math.FA], 2015; Linear Nonlinear Anal. 2 (2016), 47--52.

\bibitem{gradient}
B. Ricceri,
\emph{An alternative theorem for gradient systems},
arXiv:2002.01413v2 [math.AP], 2021.

\bibitem{misc}
B. Ricceri,
\emph{Miscellaneous Applications of Certain Minimax Theorems II},
arXiv:1810.08957 [math.FA], 2018; Acta Math. Vietnam. 45 (2020), 515--524.
\end{thebibliography}

\end{document}
